\documentclass[../portafolio.tex]{subfiles}

\begin{document}

\chapter*{About this portfolio}



This portfolio presents a collection of computational physics projects developed during the \textit{Computational Physics II} course at the Department of Physics, Faculty of Physical and Mathematical Sciences, University of Concepción.

The work showcased in this portfolio focuses on the application of scientific computing techniques to the numerical solution of problems arising in physics and applied mathematics. Through the implementation of computational algorithms in Python, these projects explore fundamental topics such as numerical differentiation, numerical integration, root-finding methods, ordinary differential equations, interpolation, curve fitting, and stochastic simulations.

A central objective of this portfolio is to demonstrate the use of computational methods as practical tools for scientific investigation. The projects emphasize the development of programming skills, algorithmic thinking, numerical modeling, data analysis, and scientific visualization, all of which are essential components of modern computational research.

Throughout these activities, several widely used scientific computing tools were employed, including Python, NumPy, Pandas, Jupyter Notebook, Git, GitHub, and \LaTeX. These technologies support reproducible research practices, version control, technical documentation, and the effective communication of scientific results.

The projects included in this portfolio reflect the application of numerical methods to a variety of physical and mathematical problems. In addition to strengthening programming proficiency, they provide practical experience in designing computational solutions, analyzing numerical accuracy, interpreting results, and communicating scientific findings through technical reports.

The skills developed through these projects form an important foundation for further work in computational physics, complex systems, data analysis, scientific modeling, and quantitative research.

\end{document}